本附录给出正文关键主张对应的source ID、来源层级和允许用途。source ID用于复核证据边界；正文中的分析推导均以这些来源为输入，但不因此变成公司披露。完整原始材料归档与来源字段由本案例证据包维护，报告不在页面中暴露本地路径。

\section{公司、重组与市场主来源}

\small
\begin{longtable}{L{1.6cm} L{4.4cm} L{2.4cm} L{6.5cm}}
\toprule
\textbf{ID} & \textbf{来源} & \textbf{层级/日期} & \textbf{主要用途与限制} \\
\midrule
\endfirsthead
\toprule
\textbf{ID} & \textbf{来源} & \textbf{层级/日期} & \textbf{主要用途与限制} \\
\midrule
\endhead
CF-00 & 600150上交所公告索引 & A1 / 截至2026-07-22 & 用于公告发现与完整性检查；索引本身不替代公告正文。 \\
CF-02 & 600150 2024年年度报告 & A1 / 2025-04-30 & 合并前legacy实际；不得作为2025年同控重述同比基数。 \\
CF-03 & 600150 2025年年度报告 & A1 / 2026-04-30 & 审计财务、2024重述比较期、聚合订单、交付、毛利、营运资金、合并日期与股本口径的核心来源；不披露逐船ASP、排程、军品经济性和统一利用率。 \\
CF-04 & 600150 2026Q1报告 & A1 / 2026-04-30 & Q1收入、归母、扣非、OCF及期末股本；未经审计，不作单季年化。 \\
CF-05 & 600150 2026H1业绩预告 & A1 / 2026-07-14 & 归母92--110亿元、扣非90--108亿元；未经审计，收入、毛利与现金未披露。 \\
CF-06--09 & 重组筹划、预案及股东批准文件 & A1 / 2024--2025 & 交易过程背景；单份文件均不证明最终实施完成。 \\
CF-10 & 上交所并购重组审核通过公告 & A1 / 2025-07-05 & 审核节点；不等同新增股份登记。 \\
CF-11 & 证监会注册批复公告 & A1 / 2025-07-19 & 注册条件；不证明全部资产权属登记或法人注销完成。 \\
CF-12 & 吸收合并报告书（最终） & A1 / 2025-07-19 & 交易条款、资产边界与交易备考；交易备考不替代法定重述实际。 \\
CF-13 & 原601989终止上市决定 & A1 / 2025-08-29 & 证明退市，不证明法人注销完成。 \\
CF-14 & 换股结果、股份变动及上市公告 & A1 / 2025-09-12 & 最终股本、新股登记与上市；公告时部分形式手续仍在办理。 \\
CF-15 & 实施情况及新增股份上市备忘录 & A1 / 2025-09-12 & 实质承继与实施；无逐资产登记完成清单。 \\
CF-16 & 2025年持续督导意见 & A1 / 2026-05-15 & 实施后监督与法律承继；形式手续缺口仍需保留。 \\
CF-17 & 2025年度暨2026Q1业绩说明会 & A1 / 2026-06-03 & 交付批次、成本占比与年度计划；只采用管理层明确答复，不采用问题前提。 \\
PR-01 & 600150盘后行情快照 & B1 / 2026-07-22 & 现价33.02元、当日高低与成交额；市值由官方股本推导，行情包0值不作为真实PE/PB。 \\
PR-02 & 前复权日K、行情、换手与流通快照 & B1 / 2026-07-22 & 727日明确\texttt{fqt=1}前复权序列、均线/回撤、量额、换手和流通A股本；供应商流通字段不得改称严格自由流通。 \\
PR-03 & 两融、沪股通、基金与龙虎榜证据包 & A1/B1 / 2026-06-30--07-22 & 上交所日度两融、季度沪股通和东方财富基金/龙虎榜查询；季度持仓不能改写为当日资金流，日频沪股通缺口有负响应。 \\
\bottomrule
\end{longtable}
\normalsize

以上来源足以支持合并估值的事实层。最重要的治理原则是：CF-03的4,674.51亿元仅为合并后上市公司民船/海工在手；CF-05仅为预告；CF-14/15证明股本与实质承继，却不允许把未披露的形式手续写成已完成。

\section{行业与政策来源}

\small
\begin{longtable}{L{1.6cm} L{4.4cm} L{2.4cm} L{6.5cm}}
\toprule
\textbf{ID} & \textbf{来源} & \textbf{层级} & \textbf{允许用途与边界} \\
\midrule
\endfirsthead
\toprule
\textbf{ID} & \textbf{来源} & \textbf{层级} & \textbf{允许用途与边界} \\
\midrule
\endhead
II-06 & 工信部2025年造船业统计 & A2 & 中国完工、新接、手持DWT与全球份额；国家数据不等于600150公司价值份额。 \\
II-07 & 中船协2025年行业运行分析 & A2 & 中国CGT份额、绿色船与成本方向；CSSC group排名不下放至600150。 \\
II-08 & Clarksons PLC 2025年报 & B1 & 全球合同额、CGT、交付与交期背景；混合美元/CGT不是船型ASP。 \\
II-09 & Clarksons绿色技术公开追踪 & B1 & 替代燃料订单与行业结构；不支持600150具体订单或利润。 \\
II-10 & OECD海事脱碳研究 & A2 & 技术、政策与全球在手背景；不证明公司订单归属。 \\
II-12 & IMO净零框架FAQ & A2 & 规则法律状态；截至截止日尚未成为生效规则，不进入基础价值。 \\
II-13--15 & UNCTAD/USTR等未完整归档线索 & 未取得/零权重 & 不支持精确数字、日期或公司影响，仅保留后续归档路径。 \\
II-16--20 & 海外同业公开材料线索 & A2/C1 & 能力与竞争背景；未完成同口径财务和估值标准化，不作可比倍数。 \\
\bottomrule
\end{longtable}
\normalsize

行业来源回答“周期处于何处、竞争壁垒是什么”，不回答“600150获得了多少订单、能赚多少利润”。因此全球TAM、中国份额、绿色化和集团排名都没有直接进入收入、EPS或目标价。

\section{券商预测与目标价证据}

\small
\begin{longtable}{L{1.7cm} L{3.0cm} L{2.2cm} L{2.4cm} L{6.0cm}}
\toprule
\textbf{ID} & \textbf{机构/材料} & \textbf{日期} & \textbf{质量} & \textbf{报告处理} \\
\midrule
\endfirsthead
\toprule
\textbf{ID} & \textbf{机构/材料} & \textbf{日期} & \textbf{质量} & \textbf{报告处理} \\
\midrule
\endhead
BR-01 & 华源H1预告点评 & 2026-07-14 & 原始PDF/B1 & 当前合并口径预测对照；目标价未披露，目标权重0。 \\
BR-02 & 诚通年报/Q1点评 & 2026-05-06 & 原始PDF/B1 & 当前合并口径预测对照；目标价未披露，目标权重0。 \\
BR-03 & 东吴年报/Q1点评 & 2026-05-05 & 原始PDF/B1 & 当前合并口径预测高端；目标价未披露，目标权重0。 \\
BR-04 & 华源Q1点评 & 2026-05-01 & 原始PDF/B1 & 与BR-01同团队，不视为独立投票；目标权重0。 \\
BR-05 & 国金Q1点评 & 2026-04-29 & 原始PDF/B1 & 当前合并口径预测对照；目标价未披露，目标权重0。 \\
BR-06--09 & 合并后历史预测 & 2025--2026 & 原始PDF/B1 & 只作预测修订史，不进入当前五份横截面。 \\
BR-10--13 & 合并前legacy报告 & 2025-08/09 & 原始PDF/B1 & 使用44.7243亿股或明确不考虑重组；当前一致预期权重0。 \\
AG-03 & 华泰58.96元第三方预览 & 2026-05-03 & 第三方预览/C2 & 未取得原始报告；不得称有效Street目标价，估值权重0。 \\
AG-04 & 国泰海通47.00元失败线索 & 2026-05-19 & 未取得/G & 本地材料无可用正文；不引用价格，估值权重0。 \\
AG-05 & iFinD一致预期失败探针 & 2026-05-18 & 未取得/G & 无可用字段，不进入报告数字。 \\
\bottomrule
\end{longtable}
\normalsize

最新五份原始报告用于2026--2028盈利预测对照，不能支持有效目标价平均。13份原始PDF全部正向但来自四家券商，评级数量不等于独立观点数量。最终目标中Broker权重为0，是证据质量约束而不是主观降权。

\section{关键主张到来源的快速映射}

\begin{exhibitbox}[复核索引]
\small
\begin{tabularx}{\exhibitboxwidth}{L{6.0cm}L{3.4cm}X}
\toprule
\textbf{关键主张} & \textbf{主source ID} & \textbf{性质} \\
\midrule
2025同控重述财务与4,674.51亿元在手 & CF-03 & 审计披露 \\
75.2562亿股、流通口径、换手与当前市值 & CF-04、CF-14、PR-01、PR-02 & 披露+机械推导 \\
复权趋势、两融、沪股通、基金与龙虎榜边界 & PR-02、PR-03 & 带日期的市场证据 \\
H1 92--110/90--108亿元 & CF-05 & 未经审计预告 \\
2026交付168艘/1,650万DWT计划 & CF-17、CF-03 & 管理计划 \\
中远约500亿元项目归属且去重 & CF-03 & 披露归属；模型增量信用0 \\
中国三大造船份额 & II-06 & 国家行业背景 \\
卖方盈利中位与无可用目标价 & BR-01--BR-13 & 原始PDF对照 \\
33.07元最终目标与31.64--35.96元区间 & 上述输入+冻结R1模型 & 分析推导 \\
\bottomrule
\end{tabularx}
\end{exhibitbox}

来源映射的投资含义是：任何新证据若要改变目标价，必须先明确它补上了哪个缺口、是否改变合并收入/毛利/归属/现金、是否与已有聚合订单重复。来源数量本身不增加估值信用。
